\documentclass[11pt]{article}
\usepackage[T1]{fontenc}
\PassOptionsToPackage{dvipsnames}{xcolor}
\usepackage{amsmath,amsthm}
\usepackage{newtxtext,newtxmath}
\usepackage[margin=1in]{geometry}
\usepackage{microtype,amsmath,amsthm,mathtools,booktabs,array,longtable}
\usepackage[dvipsnames]{xcolor}
\usepackage{enumitem,fancyhdr}
\usepackage[colorlinks=true,linkcolor=MidnightBlue,citecolor=MidnightBlue,urlcolor=MidnightBlue]{hyperref}
\usepackage{bookmark}
\hypersetup{pdftitle={Ordered Harmonic Jets: Exact Orientation Ranks, Stieltjes Commutators, and Depth-Four Hurwitz Finite Parts},pdfauthor={ProveIt research continuation},pdfsubject={Exact identities and formal orientation quotients},pdfkeywords={Hurwitz zeta, Stieltjes constants, polylogarithms, harmonic sums, Gamma function}}
\setlength{\headheight}{14pt}
\pagestyle{fancy}\fancyhf{}
\fancyhead[L]{\small Ordered harmonic jets}\fancyhead[R]{\small ProveIt research continuation}
\fancyfoot[C]{\thepage}
\setlength{\parskip}{3pt}
\setlength{\emergencystretch}{2em}
\allowdisplaybreaks[2]
\numberwithin{equation}{section}
\newtheorem{theorem}{Theorem}[section]
\newtheorem{lemma}[theorem]{Lemma}
\newtheorem{proposition}[theorem]{Proposition}
\newtheorem{corollary}[theorem]{Corollary}
\theoremstyle{definition}
\newtheorem{definition}[theorem]{Definition}
\newtheorem{example}[theorem]{Example}
\newtheorem{remark}[theorem]{Remark}
\newcommand{\R}{\mathcal R}
\newcommand{\HH}{\mathcal H}
\newcommand{\SSS}{\mathscr S}
\newcommand{\Alg}{\mathscr H}
\newcommand{\Li}{\operatorname{Li}}
\newcommand{\CT}{\operatorname{CT}}
\newcommand{\FP}{\operatorname{FP}}
\newcommand{\dd}{\,\mathrm d}
\newcommand{\Ree}{\operatorname{Re}}
\newcommand{\bbm}{\boldsymbol m}
\newcommand{\bbu}{\boldsymbol u}
\newcommand{\bbc}{\boldsymbol c}
\newcommand{\src}[1]{\texttt{\small\detokenize{#1}}}
\newcommand{\Zletter}[1]{Z_{#1}}
\title{\textbf{Ordered Harmonic Jets}\\[5pt]
\large Exact orientation ranks, Stieltjes commutators,\\
and depth-four Hurwitz finite parts}
\author{A research continuation prepared for the ProveIt project}
\date{10 October 2026}
\begin{document}
\maketitle
\begin{abstract}
The harmonic-compatible regular germs of ordered multiple Hurwitz zeta
functions leave a concrete problem: which of their mixed coefficients are
forced by one-variable identities, and which retain ordering information?
We answer this question in a precisely specified formal linear quotient.
At word length $j$ and logarithmic degree $r$, its dimension is
$\binom{j+r-1}{r}-p_{\le j}(r)$. For the complete set of coefficients entering
a depth-$d$ transverse finite part, the dimension is
$2^{d-1}-p(d)$. We give an equivariant description by augmentation spaces
of permutation modules, a rational normal-form algorithm, and a residue
argument showing that all these coordinates can be distinguished by
positive rational rays. Single-letter closure of an individual ray occurs
exactly when all its inner slopes are equal, relative to the stipulated
stuffle laws.

At depth four the quotient is the direct sum of a two-dimensional
standard permutation component and a one-dimensional sign component.
We give the full finite part and recover its three orientation coordinates
from three positive rays by an explicit rational inverse matrix. To evaluate
the coordinates analytically, we prove an all-index, absolutely convergent
Stieltjes commutator identity with its exact contact term, ordinary
polylogarithmic Mellin representations, a curvature integral, and a
Gamma-weighted antiderivative family. Rational Hurwitz shifts give explicit
cyclotomic versions. These results resolve the formal depth-four
organization question in the inspected predecessor and extend it to all
depths. They do not assert numerical period independence, reduction to a
smaller transcendental constant field, or historical priority for classical
quasi-shuffle and Hurwitz continuation principles.
\end{abstract}
\vspace{6pt}
\noindent\textbf{Keywords.} Ordered multiple Hurwitz zeta; generalized Stieltjes
constants; harmonic regularization; quasi-shuffle algebra; polylogarithm
order derivatives; Gamma-weighted primitives.
\clearpage
\tableofcontents
\clearpage

\section{The continuation target and the meaning of the result}
\label{sec:scope}
The predecessor \emph{Ordered Hurwitz Germs at Every Depth} constructs
harmonic-compatible holomorphic remainders near the all-one point. Its
second further-research question asks for the permutation components of
the lower-depth jets entering a depth-four finite part, and for a distinction
between stuffle-forced combinations and genuinely ordered coordinates
\cite{ordered}. We address that specified question, then give an all-depth
linear answer. The immediate objects are not merely numerical values: they
are words representing finite logarithmic harmonic sums, equipped with
explicit multiplication and a fixed constant-term prescription.

The distinction between three assertions will be used throughout.
An identity between evaluated functions is proved by analytic or finite
algebraic arguments. A dimension calculation in the universal word algebra
is a statement relative to its defining laws. Neither one, by itself,
proves arithmetic independence of the evaluated constants. In particular,
the number $2^{d-1}-p(d)$ below is a \emph{rational linear quotient dimension}.
It is not a transcendence degree, a count of algebraically independent
periods, or a module rank after adjoining a field of one-variable constants.

Our principal conclusions are Theorem~\ref{thm:rank}, which supplies an
all-depth normal form; Theorems~\ref{thm:identifiable} and
\ref{thm:closure}, which identify its ray coefficients and the exact
single-letter-closed locus; Theorem~\ref{thm:commutator}, which gives every
double logarithmic commutator; and Theorem~\ref{thm:quartic}, which evaluates
the full depth-four finite part in three orientation coordinates.
Theorems~\ref{thm:gamma-bridge} and \ref{thm:curvature} connect those
coordinates to Gamma-weighted integration and ordinary polylogarithmic
Mellin integrals.

\paragraph{Classical inputs and attribution.}
Finite harmonic products, their quasi-shuffle structure, and the
polynomial/Lyndon description are classical; Hoffman gives a general
algebraic framework \cite{hoffman}. Laurent expansions of multiple zeta
functions by harmonic products and analytic continuation are likewise
established subjects \cite{mow}. The analytic regularization and its
Gamma-generating diagonal used here are those of the inspected predecessor
\cite{ordered}; we reprove the required facts to make the normal-form and
identity arguments self-contained. Standard one-variable identities follow
the DLMF conventions \cite{dlmf-zeta,dlmf-polylog,dlmf-lerch,dlmf-gamma}.
No exhaustive literature-priority assertion is made for the combinations
of these tools developed here.

\paragraph{Inspection boundary.}
The manuscript entry point and selected portions of its integration,
differentiation, and research-programme chapters were read \cite{repo-manuscript,repo-discovery}, together with
the complete predecessor TeX. The incoming archives were inventoried by
visible metadata, not exhaustively unpacked or audited. The repository
reference observed during this work was
\begin{center}\small\ttfamily
d838a9eec408cb7e36f5c9bcec194314446ed295.
\end{center}
The package records source paths and hashes in its provenance manifest.
All claims of progress are relative to the explicit inspected question,
not to every unpublished file in the repository or Library.

\section{Logarithmic harmonic words and their constant terms}
\label{sec:regularization}
Fix $a>0$. Unless stated otherwise, logarithms of positive real quantities
are real logarithms. The convention for generalized Stieltjes constants is
\begin{equation}\label{eq:stieltjes}
 \zeta(1+u,a)=\frac1u+g_a(u),\qquad
 g_a(u)=\sum_{m\ge0}\frac{(-1)^m\gamma_m(a)}{m!}u^m.
\end{equation}
In particular $\gamma_0(a)=-\psi(a)$. A superscript on $\zeta^{(m)}$
always means differentiation in the spectral argument; a derivative in
$a$ will be written explicitly.

A letter is a pair $(k,m)$ with $k\ge1$ and $m\ge0$. It represents
$f_{k,m}(x)=x^{-k}\log^m x$. For a word
$w=((k_1,m_1),\ldots,(k_j,m_j))$, put
\begin{equation}\label{eq:finite-word}
 W_w(N;a)=\sum_{N>n_1>\cdots>n_j\ge0}
             \prod_{i=1}^j f_{k_i,m_i}(n_i+a),\qquad W_\varnothing=1.
\end{equation}
The larger index carries the earlier letter. The finite generalized
harmonic numbers used below are
\begin{equation}\label{eq:harmonic}
 H_N^{(k,m)}(a)=\sum_{n=0}^{N-1}\frac{\log^m(n+a)}{(n+a)^k}.
\end{equation}
Thus $H_N^{(1,0)}(1)$ is the ordinary $N$th harmonic number.

\begin{proposition}[Harmonic constant-term character]\label{prop:character}
Every word has a unique polynomial $P_w(L;a)$ such that, with
$L_N=\log(N+a)$,
\begin{equation}\label{eq:polynomial-cutoff}
 W_w(N;a)=P_w(L_N;a)+O\bigl(N^{-1}(1+\log N)^M\bigr)
\end{equation}
for some $M$. Define $\R_a(w)=P_w(0;a)$. Let $*$ be the positive-merge
quasi-shuffle product, with
\begin{equation}\label{eq:merge}
 (k,m)\circ(\ell,n)=(k+\ell,m+n).
\end{equation}
Then
\begin{equation}\label{eq:character}
 \R_a(w*v)=\R_a(w)\R_a(v).
\end{equation}
It agrees with every convergent sum in this class, and
\begin{equation}\label{eq:single-letters}
 \ell_{k,m}(a):=\R_a((k,m))=
 \begin{cases}
 \gamma_m(a),&k=1,\\
 (-1)^m\zeta^{(m)}(k,a),&k\ge2.
 \end{cases}
\end{equation}
The estimates, and all fixed parameter derivatives needed below, are
locally uniform for $a$ in compact subsets of $(0,\infty)$.
\end{proposition}
\begin{proof}
Induct on length. Finite inner sums have at most logarithmic growth. If
the first exponent is at least two, the outer series converges and its
tail has the form of the error in \eqref{eq:polynomial-cutoff}. If the
first letter is $(1,m)$, insert the tail polynomial into the exact increment
\[
 W_w(N+1;a)-W_w(N;a)=
       \frac{\log^m(N+a)}{N+a}W_{w'}(N;a).
\]
The primitive in $L$ of $L^mP_{w'}(L;a)$ is a polynomial. The sum--integral
difference of this increment is summable, since its derivative in the
continuous variable is $O(x^{-2}(1+\log x)^{M'})$. The inserted tail error
is summable for the same reason. Its accumulated limit defines the missing
constant. This proves existence and the error estimate; uniqueness follows
because no nonzero polynomial in $\log(N+a)$ tends to zero.

At finite cutoff, interlacing the two sets of indices and accounting for
equalities proves the quasi-shuffle identity. Products of the error terms
with a fixed logarithmic polynomial still tend to zero. Hence the
polynomial asymptotics satisfy the same product identity. Evaluation at
$L=0$ is multiplicative. The depth-one identities follow by Taylor
expanding the Hurwitz tail, or directly from \eqref{eq:stieltjes} and the
convergent spectral derivatives at $k\ge2$. The increment arguments and
summable majorants remain uniform on compact positive $a$-intervals.
Differentiation of those majorants gives the stated parameter regularity.
\end{proof}

We write $Z_m=(1,m)$, and frequently $X=Z_0$, $Y=Z_1$, $Z=Z_2$.
Concatenation is not the product $*$: for instance
\begin{equation}\label{eq:double-stuffle}
 \R_a(Z_pZ_q)+\R_a(Z_qZ_p)
 =\gamma_p(a)\gamma_q(a)-(-1)^{p+q}\zeta^{(p+q)}(2,a).
\end{equation}
This elementary distinction is essential at higher depth.

\begin{proposition}[Shift and differential transport]\label{prop:transport}
For a nonempty word $w=v(k,m)$,
\begin{equation}\label{eq:word-shift}
 \R_a(w)-\R_{a+1}(w)=\frac{\log^m a}{a^k}\R_{a+1}(v).
\end{equation}
Define a derivation on concatenated words by the letter rule
\begin{equation}\label{eq:letter-D}
 \mathcal D(k,m)=-k(k+1,m)+m(k+1,m-1),
\end{equation}
where a term with zero coefficient is omitted. Then
\begin{equation}\label{eq:word-D}
 \frac{\mathrm d}{\mathrm da}\R_a(w)=\R_a(\mathcal Dw).
\end{equation}
\end{proposition}
\begin{proof}
Separate the finite sum according to whether its smallest index is zero.
The sum with all indices positive becomes a sum at $a+1$ after subtracting
one from every index. The two cutoff logarithms can be chosen identical:
$\log(N+a)=\log((N-1)+(a+1))$. Constant terms give
\eqref{eq:word-shift}. Differentiating each summand gives
\eqref{eq:letter-D}. The locally uniform asymptotic construction in
Proposition~\ref{prop:character} permits differentiation of the polynomial
asymptotic. Derivatives of $L_N$ are $O(N^{-1})$, so contribute no constant
term. This proves \eqref{eq:word-D}.
\end{proof}

\section{Compatible ordered germs and transverse finite parts}
\label{sec:germs}
Define
\[
 Z_d(s_1,\ldots,s_d;a)=\sum_{n_1>\cdots>n_d\ge0}
                         \prod_{i=1}^d(n_i+a)^{-s_i},\qquad Z_0=1.
\]
Initially all $\Ree s_i>1$ suffice. The following construction provides
the only local continuation needed here.

Put
\begin{equation}\label{eq:kernel-B}
 K(t)=\frac1{e^t-1}-\frac1t,\qquad
 B(u,x)=\zeta(1+u,x+1)-\frac{x^{-u}}u.
\end{equation}
The singularity at $u=0$ is removable, and
\begin{equation}\label{eq:B-Mellin}
 B(u,x)=\frac1{\Gamma(1+u)}\int_0^\infty t^u e^{-xt}K(t)\dd t,
 \qquad \Ree u>-1.
\end{equation}
The integral follows by subtracting the usual Hurwitz and elementary
Mellin integrals where they converge, then continuing. Since
$K(t)=-1/2+O(t)$ at zero, scaling $t=y/x$ shows that $B$ and its fixed
$u$-derivatives are $O(x^{-1+\delta}(1+\log x)^M)$ on small compact
$u$-sets. In particular,
\begin{equation}\label{eq:tail-A}
 A_d(\bbu;a)=\sum_{n_2>\cdots>n_d\ge0}
 B(u_1,n_2+a)\prod_{i=2}^d(n_i+a)^{-1-u_i}
\end{equation}
converges normally in a sufficiently small neighborhood of zero, including
after any fixed mixed derivative. Grouping by $n_2$ gives a majorant
$O(n_2^{-2+\delta}(1+\log n_2)^M)$, with $\delta<1$.

\begin{proposition}[Compatible germ and harmonic jets]\label{prop:germ}
Set $\HH_0=1$, $\HH_1(u;a)=g_a(u)$, and, recursively,
\begin{equation}\label{eq:H-recursion}
 \HH_d(\bbu;a)=A_d(\bbu;a)+
 \frac{\HH_{d-1}(u_1+u_2,u_3,\ldots;a)
       -\HH_{d-1}(u_2,u_3,\ldots;a)}{u_1}.
\end{equation}
The quotient is removable. These functions are holomorphic near zero and,
with $U_k=u_1+\cdots+u_k$,
\begin{equation}\label{eq:germ}
 Z_d(1+u_1,\ldots,1+u_d;a)=
 \sum_{j=0}^d\frac{\HH_j(u_{d-j+1},\ldots,u_d;a)}{U_1\cdots U_{d-j}}.
\end{equation}
Their Taylor coefficients are
\begin{equation}\label{eq:jet-word}
 [u_1^{m_1}\cdots u_j^{m_j}]\HH_j
 =\frac{(-1)^{m_1+\cdots+m_j}}{m_1!\cdots m_j!}
                       \R_a(Z_{m_1}\cdots Z_{m_j}).
\end{equation}
\end{proposition}
\begin{proof}
Sum the outer index first and use \eqref{eq:kernel-B}. This gives
\[
 Z_d(1+\bbu;a)=\frac1{u_1}
 Z_{d-1}(1+u_1+u_2,1+u_3,\ldots;a)+A_d(\bbu;a).
\]
Insert the lower-depth instance of \eqref{eq:germ}, then add and subtract
$\HH_{d-1}(u_2,\ldots)/u_1$. Induction proves the formula and
holomorphy, since a divided difference of a holomorphic function is
holomorphic.

To identify the cutoff prescription rather than merely a formal
subtraction, construct locally holomorphic functions $Q_j(\bbu;L)$ with
\[
 Q_j(\bbu;L)=k_j(\bbu)+\int_0^L e^{-u_1t}
                             Q_{j-1}(u_2,\ldots;t)\dd t.
\]
The increment proof of Proposition~\ref{prop:character}, now with small
complex exponent perturbations, shows that the finite ordered sums equal
$Q_j(\bbu;\log(N+a))+O(N^{-1+\delta}(1+\log N)^M)$ locally uniformly.
Indeed, the difference between a discrete increment and the integral over
$[N+a,N+a+1]$ is $O(N^{-2+\delta}(1+\log N)^M)$; so is the error from
inserting the lower-depth approximation. Their normally convergent sum
constructs $k_j$. Cauchy estimates on a slightly larger compact set justify
all fixed Taylor extractions.

When all $\Ree u_i>0$, the finite sums converge. Repeated integration of
the displayed recurrence gives the prefix denominators in
\eqref{eq:germ}. Uniqueness by subtraction at successive depths therefore
identifies $k_j=\HH_j$. Taylor extraction produces a polynomial in $L$
whose constant is the coefficient of $\HH_j$. Extracting the same
coefficient from the finite sum proves \eqref{eq:jet-word}.
\end{proof}

A transverse ray for the present purpose is a vector
$\bbc=(c_1,\ldots,c_d)$ with every prefix sum
$S_k=c_1+\cdots+c_k\ne0$. Its finite part is the constant Laurent
coefficient
\[
 F_d(\bbc;a)=\FP_{\varepsilon=0}
       Z_d(1+c_1\varepsilon,\ldots,1+c_d\varepsilon;a).
\]
The vector may have mixed signs; positive vectors are automatically
admissible. No choice of approaching half-plane is needed for the local
meromorphic coefficient.

\begin{corollary}[Exact ray coefficient formula]\label{cor:ray}
For every admissible ray,
\begin{equation}\label{eq:ray}
 F_d(\bbc;a)=\sum_{j=1}^d
 \frac{(-1)^{d-j}}{\prod_{k=1}^{d-j}S_k}
 \sum_{\substack{m_1+\cdots+m_j=d-j\\m_i\ge0}}
 \frac{\prod_{i=1}^j c_{d-j+i}^{m_i}}{\prod_{i=1}^j m_i!}
                      \R_a(Z_{m_1}\cdots Z_{m_j}).
\end{equation}
\end{corollary}
\begin{proof}
Each denominator in \eqref{eq:germ} is homogeneous of degree $d-j$ in
$\varepsilon$. Extract exactly that Taylor degree from its numerator and
use \eqref{eq:jet-word}. The $j=0$ term is a pure pole and has no constant
coefficient.
\end{proof}

\section{Exact formal ranks and permutation components}
\label{sec:ranks}
Let $\Alg$ be the rational vector space freely spanned by all decorated
words, with the quasi-shuffle product \eqref{eq:merge}. Its first weight
is $\sum k_i$ and its logarithmic degree is $\sum m_i$. Let $\SSS$ be
its unital subalgebra generated by all individual letters. This is the
\emph{single-letter algebra}: evaluation sends it to polynomials in
ordinary Stieltjes constants and spectral derivatives of Hurwitz zeta.

For fixed $j\ge1,r\ge0$, let $V_{j,r}$ be the span of the words
$Z_{m_1}\cdots Z_{m_j}$ with $\sum m_i=r$. We ask for its linear quotient
modulo the elements already in $\SSS$. This completely specifies the
meaning of ``orientation'' in the rank theorem.

\begin{lemma}[Finite symmetric-sum identity]\label{lem:sym-sum}
For arbitrary $m_1,\ldots,m_j\ge0$, including repeated values,
\begin{equation}\label{eq:symmetric-sum}
 \sum_{\sigma\in S_j}\R_a(Z_{m_{\sigma(1)}}\cdots Z_{m_{\sigma(j)}})
 =\sum_{\pi\in\Pi_j}(-1)^{j-|\pi|}
       \prod_{B\in\pi} (|B|-1)!\,
                  \ell_{|B|,\sum_{i\in B}m_i}(a).
\end{equation}
The same identity holds in $\Alg$ before evaluation, with ordinary
products on the right replaced by quasi-shuffle products. The left sum is
over labelled permutations, not only distinct words.
\end{lemma}
\begin{proof}
At finite cutoff, sum over $j$ labelled indices with the condition that
all are distinct. Inclusion--exclusion on their equality partitions has
partition-lattice coefficient
$\prod_B(-1)^{|B|-1}(|B|-1)!$. This coefficient can also be verified
inductively: the sum of signed cycle permutations on a set of size $n$
is zero for $n>1$, and one for $n=1$. A block $B$ of equal indices gives
the single letter $(|B|,\sum_{i\in B}m_i)$. Ordering the distinct indices
gives the labelled permutation sum. Taking constant terms proves
\eqref{eq:symmetric-sum}. The same combinatorial argument interlaces
formal words and proves the pre-evaluation identity.
\end{proof}

\begin{theorem}[Orientation quotient and all-depth count]\label{thm:rank}
The intersection $V_{j,r}\cap\SSS$ is exactly the span of distinct-word
permutation-orbit sums. Consequently,
\begin{equation}\label{eq:rank-jr}
 \dim_{\mathbb Q}\frac{V_{j,r}}{V_{j,r}\cap\SSS}
       =\binom{j+r-1}{r}-p_{\le j}(r),
\end{equation}
where $p_{\le j}(r)$ counts partitions of $r$ into at most $j$ parts,
with $p_{\le j}(0)=1$. In the direct sum of all spaces entering
\eqref{eq:ray}, the quotient dimension is
\begin{equation}\label{eq:rank-d}
 \boxed{\displaystyle
 N_d=\dim_{\mathbb Q}
 \frac{\bigoplus_{j=1}^d V_{j,d-j}}
      {\bigl(\bigoplus_{j=1}^d V_{j,d-j}\bigr)\cap\SSS}
      =2^{d-1}-p(d).}
\end{equation}
\end{theorem}
\begin{proof}
A quasi-shuffle product of $h$ single letters has word length at most
$h$. Its length-$h$ component is the full shuffle symmetrization of those
letters. A monomial in single letters of first weight $j$ has at most
$j$ factors. To have a component of length $j$, it must have exactly
$j$ factors, each with first exponent one. Thus the length-$j$
component of every element of $\SSS$ of bidegree $(j,r)$ is a linear
combination of orbit sums. If that element belongs to $V_{j,r}$, there
are no other lengths and the necessity follows. Conversely, every orbit
sum belongs to $\SSS$ by Lemma~\ref{lem:sym-sum}, after dividing the
labelled sum by the product of the multiplicity factorials.

The number of raw words is the weak-composition count
$\binom{j+r-1}{r}$. Their orbits are the nondecreasing $j$-tuples of sum
$r$, counted by $p_{\le j}(r)$. Distinct orbit sums have disjoint word
supports, so are linearly independent. This proves \eqref{eq:rank-jr}.

The spaces with bidegrees $(j,d-j)$ cannot mix, since $\SSS$ is
bigraded. There are $\sum_j\binom{d-1}{j-1}=2^{d-1}$ raw words.
Adding one to every entry of a nondecreasing tuple $(m_1,\ldots,m_j)$
of sum $d-j$ gives a partition of $d$ into exactly $j$ positive parts.
Summing the orbit counts over $j$ therefore gives $p(d)$, proving
\eqref{eq:rank-d}.
\end{proof}

\begin{corollary}[Equivariant description]\label{cor:representation}
For each nondecreasing tuple $\bbm$ of length $j$ and sum $r$, let
$G_{\bbm}\le S_j$ be the subgroup permuting equal entries. As an
$S_j$-representation,
\begin{equation}\label{eq:permutation-modules}
 \frac{V_{j,r}}{V_{j,r}\cap\SSS}
 \simeq\bigoplus_{\bbm}
 \left(\mathbb Q[S_j/G_{\bbm}]/\mathbb Q\mathbf1_{\bbm}\right).
\end{equation}
Equivalently each summand is the sum-zero subspace on that orbit.
\end{corollary}
\begin{proof}
Each orbit is its permutation representation on the cosets of its
stabilizer. Theorem~\ref{thm:rank} quotients out exactly its invariant
orbit-sum vector. Averaging its coefficients splits off that line over
$\mathbb Q$, leaving the sum-zero subspace.
\end{proof}

Here is a direct rational normal-form algorithm. In each orbit choose a
canonical word $w_0$, for example the nondecreasing tuple. For every other
word keep $\delta_w=w-w_0$. If the orbit has size $h$ and its known
orbit sum is $T_{\mathcal O}$, then
\begin{equation}\label{eq:normal-algorithm}
 w_0=\frac1h\left(T_{\mathcal O}-\sum_{w\ne w_0}\delta_w\right),
 \qquad w=w_0+\delta_w.
\end{equation}
Lemma~\ref{lem:sym-sum} evaluates $T_{\mathcal O}$ explicitly in single
letters. Formula \eqref{eq:normal-algorithm} is a terminating finite
certificate, not a numerical relation search.

\begin{center}
\begin{tabular}{rrrr}
\toprule
Depth $d$ & Raw coefficients $2^{d-1}$ & Symmetric orbits $p(d)$ & $N_d$\\
\midrule
1&1&1&0\\2&2&2&0\\3&4&3&1\\4&8&5&3\\
5&16&7&9\\6&32&11&21\\7&64&15&49\\8&128&22&106\\
9&256&30&226\\10&512&42&470\\11&1024&56&968\\12&2048&77&1971\\
\bottomrule
\end{tabular}
\end{center}

\begin{remark}[What is not counted]\label{rem:not-module-rank}
The theorem does not count indecomposable Lyndon generators. Products
containing an already ordered constant may relate one orientation
coordinate to another and to a lower-depth ordered quantity. Nor does it
compute a quotient after making all single-letter values scalars in a
larger coefficient field. The first-weight/log-degree grading and the
specified \emph{rational linear} quotient are indispensable. Section
\ref{sec:quartic} exhibits this distinction explicitly: a triple
orientation coordinate is related to $\gamma_0\Omega$ and to a merged-letter
double sum, without ceasing to be a nonzero class in the quotient just
computed.
\end{remark}

\section{Ray identifiability and the exact single-letter-closed locus}
\label{sec:ray-geometry}
The quotient count is useful only if the coefficients of the rays can
actually distinguish the coordinates. The nested polar denominators make
this possible without numerical sampling assumptions.

\begin{theorem}[All raw jet coefficients are identifiable]\label{thm:identifiable}
For fixed $d$, the rational functions
\begin{equation}\label{eq:coefficient-functions}
 f_{j,\bbm}(\bbc)=
 \frac{\prod_{i=1}^j c_{d-j+i}^{m_i}}
      {\prod_{k=1}^{d-j}S_k\,\prod_i m_i!},
 \qquad |\bbm|=d-j,
\end{equation}
are linearly independent over $\mathbb Q$. Consequently the $N_d$
orientation coefficient functions in any rational normal form are also
linearly independent. Finitely many positive rational rays distinguish
all of them; $N_d$ such rays can be chosen for the orientation quotient.
\end{theorem}
\begin{proof}
Suppose a rational linear combination of \eqref{eq:coefficient-functions}
vanishes. Collect its terms at a fixed $j$ into a homogeneous polynomial
$P_j(c_{d-j+1},\ldots,c_d)$ over the denominator
$S_1\cdots S_{d-j}$. Restrict the residue first to $S_{d-1}=0$, away
from the other prefix hyperplanes. Only the $j=1$ term can contribute,
so $P_1(c_d)=0$. The suffix variable $c_d$ is unconstrained on that
hyperplane, hence $P_1$ is the zero polynomial.

Next take the residue at $S_{d-2}=0$. After the $j=1$ term is removed,
only $j=2$ contributes. Its two suffix variables remain arbitrary, so
$P_2=0$. Continue in this fashion. The last term has no denominator and
must also vanish. Within each suffix the distinct monomials are linearly
independent, proving the first assertion.

A change from raw words to orbit sums and orientation differences is an
invertible rational change of coordinates, so its full set of coefficient
functions remains independent. Every subset, including the orientation
functions, is independent. If their values on all positive rational rays
spanned a proper subspace, a nonzero linear functional would annihilate
all of them there. Density and continuity would give an identity on the
positive open orthant, hence a rational-function identity, a contradiction.
Successively choose rays whose value vectors increase the span.
\end{proof}

In this statement ``distinguish'' means that, after the single-letter
contribution is subtracted, an exact finite linear system recovers the
formal orientation values. It does not imply that the recovered numerical
values are independent of one another or of other periods.

\begin{theorem}[Complete formal single-letter closure]\label{thm:closure}
Let $d\ge3$, and let every prefix sum of $\bbc$ be nonzero. Replace the
values in \eqref{eq:ray} by their universal words. That formal finite part
belongs to the single-letter algebra (with the slope coefficients adjoined)
if and only if
\begin{equation}\label{eq:closed-locus}
 c_2=c_3=\cdots=c_d.
\end{equation}
\end{theorem}
\begin{proof}
The component of bidegree $(d-1,1)$ in \eqref{eq:ray} is
\[
 -\frac1{c_1}\sum_{i=1}^{d-1}c_{i+1}
                   X^{i-1}YX^{d-1-i}.
\]
The words here constitute a single permutation orbit. By Theorem
\ref{thm:rank}, its coefficient vector belongs to $\SSS$ exactly when
all its entries are equal. This forces \eqref{eq:closed-locus}.
Conversely, if every inner slope is $t$, each suffix polynomial in
\eqref{eq:ray} is invariant under permuting its letters: the common factor
is $t^{d-j}$ and the factorial denominator is permutation invariant.
Every such combination is an orbit-sum combination, and hence lies in
$\SSS$. The distinct bidegrees ensure that neither direction can be
altered by cancellation from another layer.
\end{proof}

\begin{remark}
This is a necessary-and-sufficient result within a fixed universal law
system. At particular numerical $a$, an extra analytic identity could
reduce a ray outside \eqref{eq:closed-locus}. The theorem does not rule
out that possibility. At $d=2$ there is no nontrivial orientation quotient,
and the one-entry inner-slope condition is vacuous.
\end{remark}

\section{Gamma generation of every closed ray}
\label{sec:gamma-generation}
The predecessor's diagonal generator provides an explicit evaluation on
the complete closed locus just found. Define, as a formal power series in
$z$ with coefficients holomorphic near $u=0$,
\begin{equation}\label{eq:master-E}
 \mathcal E_a(z,u)=\exp\left\{
 zg_a(u)+\sum_{k\ge2}\frac{(-1)^{k+1}}k\zeta(k+ku,a)z^k\right\}.
\end{equation}

\begin{proposition}[Diagonal generator]\label{prop:diagonal-generator}
For every $j\ge0$,
\begin{equation}\label{eq:E-H}
 [z^j]\mathcal E_a(z,u)=\HH_j(u,\ldots,u;a),\qquad
 \mathcal E_a(z,0)=\frac{\Gamma(a)}{\Gamma(a+z)}.
\end{equation}
\end{proposition}
\begin{proof}
For positive $u$, the elementary-symmetric generating identity for the
convergent diagonal multiple zeta sums is
\[
 \sum_{d\ge0}Z_d(1+u,\ldots,1+u;a)z^d
 =\exp\left\{\sum_{k\ge1}\frac{(-1)^{k+1}}k\zeta(k+ku,a)z^k\right\}.
\]
The prefix denominators in \eqref{eq:germ} on this diagonal are
$(d-j)!u^{d-j}$. Therefore its left side equals
$e^{z/u}\sum_{j\ge0}\HH_j(u,\ldots,u;a)z^j$. Cancel $e^{z/u}$
coefficientwise and continue the holomorphic coefficients to $u=0$.
The Taylor expansion of $\log\Gamma(a+z)$ then proves the second
identity, since
\[
 \log\frac{\Gamma(a)}{\Gamma(a+z)}
 =\gamma_0(a)z+\sum_{k\ge2}\frac{(-1)^{k+1}}k\zeta(k,a)z^k.
\]
\end{proof}

\begin{corollary}[All-depth closed-ray identity]\label{cor:closed-ray}
When $c+rt\ne0$ for $0\le r\le d-1$,
\begin{equation}\label{eq:closed-ray}
 \boxed{\displaystyle
 F_d(c,t,\ldots,t;a)=\sum_{j=1}^d
 \frac{t^{d-j}}{\prod_{r=0}^{d-j-1}(c+rt)}
                    [z^ju^{d-j}]\mathcal E_a(z,u).}
\end{equation}
Every coefficient on the right is a finite polynomial in generalized
Stieltjes constants and spectral derivatives $\zeta^{(m)}(k,a)$,
$k\ge2$. At $t=0$ this gives $F_d(c,0,\ldots,0;a)=C_d(a)$, where
\begin{equation}\label{eq:C-gen}
 C_d(a)=[z^d]\frac{\Gamma(a)}{\Gamma(a+z)}=\R_a(X^d).
\end{equation}
\end{corollary}
\begin{proof}
In \eqref{eq:germ}, the $j$th numerator along this ray is
$\HH_j(t\varepsilon,\ldots,t\varepsilon;a)$. Its denominator is
$\varepsilon^{d-j}\prod_{r=0}^{d-j-1}(c+rt)$. Use
\eqref{eq:E-H} and extract the constant coefficient. The $t=0$ value
follows directly, or by continuity. Finite coefficient extraction from
\eqref{eq:master-E} uses only finitely many $k$ and finitely many spectral
derivatives.
\end{proof}

For later use, put $x=\gamma_0(a)$ and $z_k=\zeta(k,a)$. Then
\begin{align}
 C_2&=\frac{x^2-z_2}{2},\qquad
 C_3=\frac{x^3-3xz_2+2z_3}{6},\label{eq:C23}\\
 C_4&=\frac{x^4-6x^2z_2+3z_2^2+8xz_3-6z_4}{24}.
 \label{eq:C4}
\end{align}
Because $z_k=(-1)^k\psi^{(k-1)}(a)/(k-1)!$ for $k\ge2$,
the $C_j$ themselves are polygamma polynomials.

\section{Every double orientation as a convergent Stieltjes identity}
\label{sec:commutators}
Define the logarithmic commutators
\begin{equation}\label{eq:commutator-def}
 \mathcal C_{p,q}(a)=\R_a(Z_pZ_q)-\R_a(Z_qZ_p),\qquad p,q\ge0.
\end{equation}
They are antisymmetric in $p,q$. Combined with
\eqref{eq:double-stuffle}, one commutator gives both ordered double
values. The all-index contact term in the next theorem is part of the
identity, not a matter of notation.

\begin{theorem}[All-index Stieltjes commutator]\label{thm:commutator}
Let $p,q\ge0$,
\[
 \kappa_{p,q}=\frac1{q+1}-\frac1{p+1},\qquad n=p+q+1.
\]
Then, for every $a>0$,
\begin{equation}\label{eq:commutator-series}
 \boxed{\displaystyle
 \mathcal C_{p,q}(a)=\kappa_{p,q}\gamma_n(a)+
 \sum_{\nu=0}^\infty
 \frac{\log^q x\,\gamma_p(x)-\log^p x\,\gamma_q(x)
                  -\kappa_{p,q}\log^n x}{x},\quad x=\nu+a.}
\end{equation}
The series is ordinarily absolutely convergent, locally uniformly in $a$.
For $p\ne q$, its summand has the expansion
\begin{equation}\label{eq:commutator-tail}
 \frac{q-p}{12x^3}\log^{p+q-1}x
       +O\!\left(\frac{1+\log^{p+q-1}x}{x^5}\right).
\end{equation}
For $p=q$ both sides vanish identically.
\end{theorem}
\begin{proof}
At depth two, \eqref{eq:H-recursion} reads
\[
 \HH_2(u,v)=A_2(u,v)+\frac{g_a(u+v)-g_a(v)}u.
\]
The contribution of the divided difference to the coefficient difference
corresponding to \eqref{eq:commutator-def} is
\begin{align*}
 &(-1)^{p+q}p!q!\,[u^pv^q-u^qv^p]
                   \frac{g_a(u+v)-g_a(v)}u\\
 &\quad=-\frac{p!q!}{n!}
     \left\{\binom n{p+1}-\binom n{q+1}\right\}\gamma_n(a)
       =\kappa_{p,q}\gamma_n(a).
\end{align*}
Here bracket notation denotes the difference of the two indicated
coefficient extractions. This computation fixes the contact term exactly.

For $x>0$, Taylor expansion of \eqref{eq:kernel-B} gives
\begin{equation}\label{eq:B-derivative}
 \partial_u^mB(0,x)=(-1)^m
       \left\{\gamma_m(x+1)+\frac{\log^{m+1}x}{m+1}\right\}.
\end{equation}
Differentiate the normally convergent series
$A_2(u,v)=\sum_{\nu\ge0}B(u,\nu+a)(\nu+a)^{-1-v}$.
Using \eqref{eq:B-derivative}, its antisymmetric contribution is
\[
 \sum_{\nu\ge0}\frac1x
 \left\{\log^q x\,\gamma_p(x+1)
       -\log^p x\,\gamma_q(x+1)
       -\kappa_{p,q}\log^{p+q+1}x\right\}.
\]
The Stieltjes shift
$\gamma_m(x+1)=\gamma_m(x)-\log^m x/x$ cancels between the two
terms, proving \eqref{eq:commutator-series}.

For completeness, the first Euler--Maclaurin terms of the one-variable
Hurwitz function give
\begin{equation}\label{eq:gamma-asym}
 \gamma_m(x)=-\frac{L^{m+1}}{m+1}+\frac{L^m}{2x}
       +\frac{L^m-mL^{m-1}}{12x^2}
       +O\bigl(x^{-4}(1+L^m)\bigr),\quad L=\log x.
\end{equation}
The term $mL^{m-1}$ is zero when $m=0$. To obtain the sharper
antisymmetric remainder in \eqref{eq:commutator-tail}, retain the next
Euler--Maclaurin term before taking the difference. It is $x^{-4}$ times
a polynomial obtained from $(1+u)_3e^{-uL}$; its top $L^m$ coefficient is
independent of $m$, so the top $L^{p+q}$ term cancels. The remaining degree
is at most $p+q-1$. The same differentiated remainder argument, or one
additional Euler--Maclaurin term, proves the stated $O$ estimate. The
$x^0$ contribution in the numerator of
\eqref{eq:commutator-series} is exactly cancelled by the displayed
logarithmic subtraction; the $x^{-1}$ terms cancel each other. The
$x^{-2}$ term is $(q-p)L^{p+q-1}/12$. Division by $x$ proves
\eqref{eq:commutator-tail} and absolute convergence.
\end{proof}

\begin{corollary}[Explicit ordered double normal form]\label{cor:double-normal}
For every $p,q\ge0$,
\begin{equation}\label{eq:double-normal}
 \R_a(Z_pZ_q)=\frac12\left\{
 \gamma_p(a)\gamma_q(a)-(-1)^{p+q}\zeta^{(p+q)}(2,a)
                         +\mathcal C_{p,q}(a)\right\}.
\end{equation}
Moreover,
\begin{equation}\label{eq:C-shift}
 \mathcal C_{p,q}(a)-\mathcal C_{p,q}(a+1)
   =\frac{\log^q a\,\gamma_p(a)-\log^p a\,\gamma_q(a)}a.
\end{equation}
\end{corollary}
\begin{proof}
Add and subtract \eqref{eq:double-stuffle} and
\eqref{eq:commutator-def}. Apply \eqref{eq:word-shift} to the two
ordered words for the second identity; the products of the two
$a$-terms cancel.
\end{proof}

Write $\Omega=\mathcal C_{0,1}$ and
$\Omega_2=\mathcal C_{0,2}$. The first is precisely the orientation
constant in the inspected predecessor. The theorem reproduces its
normalization, including the $-\gamma_2/2$ term, and gives the next one:
\begin{align}
 \Omega(a)&=-\frac{\gamma_2(a)}2+
 \sum_{\nu\ge0}\frac{L\gamma_0(x)-\gamma_1(x)+L^2/2}{x},
 \label{eq:Omega}\\
 \Omega_2(a)&=-\frac{2\gamma_3(a)}3+
 \sum_{\nu\ge0}\frac{L^2\gamma_0(x)-\gamma_2(x)+2L^3/3}{x},
 \label{eq:Omega2}
\end{align}
where $x=\nu+a$ and $L=\log x$. Thus these are exact analytic
coordinates, rather than residuals of an integer-relation search.

\section{Ordinary polylogarithmic Mellin identities}
\label{sec:mellin}
For $a,t>0$, define
\begin{equation}\label{eq:Lerch-heat}
 \mathcal L_a(s;t)=\sum_{\nu\ge0}\frac{e^{-(\nu+a)t}}{(\nu+a)^s}
                 =e^{-at}\Phi(e^{-t},s,a).
\end{equation}
This is entire in $s$ for each fixed $t>0$. At $a=1$ it is
$\Li_s(e^{-t})$. Define the elementary reciprocal-Gamma polynomials
\begin{equation}\label{eq:P-polynomials}
 P_m(T)=\left.\partial_u^m\frac{e^{uT}}{\Gamma(1+u)}\right|_{u=0}.
\end{equation}
With $\gamma$ denoting Euler's constant, the first four are
\begin{align*}
 P_0&=1,&P_1&=T+\gamma,\\
 P_2&=(T+\gamma)^2-\zeta(2),&
 P_3&=(T+\gamma)^3-3\zeta(2)(T+\gamma)+2\zeta(3).
\end{align*}
All higher ones follow by finite coefficient extraction from the
Taylor expansion of $1/\Gamma(1+u)$.

\begin{theorem}[Mellin realization of every commutator]\label{thm:C-Mellin}
For $p,q\ge0$, with $n$ and $\kappa_{p,q}$ as in
Theorem~\ref{thm:commutator},
\begin{align}\label{eq:C-Mellin}
 \mathcal C_{p,q}(a)-\kappa_{p,q}\gamma_n(a)
 ={}&(-1)^{p+q}\int_0^\infty K(t)
 \bigl\{P_p(\log t)\,\partial_s^q\mathcal L_a(1;t)\notag\\
 &\hspace{39mm}-P_q(\log t)\,\partial_s^p\mathcal L_a(1;t)\bigr\}\dd t.
\end{align}
This is an ordinary absolutely convergent integral; no finite-part
interpretation is needed at either endpoint.
\end{theorem}
\begin{proof}
Insert \eqref{eq:B-Mellin} into $A_2$ and interchange a normally
convergent sum and integral to obtain
\[
 A_2(u,v)=\frac1{\Gamma(1+u)}\int_0^\infty
                       K(t)t^u\mathcal L_a(1+v;t)\dd t.
\]
Differentiate in $u,v$ and antisymmetrize, using
\eqref{eq:jet-word} and the already computed divided-difference term.
For a direct endpoint check, $\partial_s^m\mathcal L_a(1;t)$ grows at
most like a fixed power of $1+|\log t|$ as $t\downarrow0$. This follows
by splitting its defining sum at $\nu\asymp1/t$ and using the exponential
tail. The factor $K(t)$ is bounded there. At infinity every order
derivative decays exponentially. These estimates justify all the fixed
differentiations and prove the asserted convergence.
\end{proof}

At $a=1$, \eqref{eq:C-Mellin} supplies an all-index family involving
ordinary polylogarithms and their order derivatives. For example,
\begin{align}\label{eq:Omega-Mellin}
 \Omega(1)+\frac{\gamma_2}{2}
 =\int_0^\infty K(t)\biggl\{
 (\log t+\gamma)\Li_1(e^{-t})
 -\left.\partial_s\Li_s(e^{-t})\right|_{s=1}\biggr\}\dd t,\\
 \Omega_2(1)+\frac{2\gamma_3}{3}
 =\int_0^\infty K(t)\biggl\{
 \left.\partial_s^2\Li_s(e^{-t})\right|_{s=1}
 -\bigl((\log t+\gamma)^2-\zeta(2)\bigr)\Li_1(e^{-t})
 \biggr\}\dd t.
 \label{eq:Omega2-Mellin}
\end{align}
Here $\Li_1(e^{-t})=-\log(1-e^{-t})$. No relation between
$\Omega(1)$ and a smaller specified constant algebra is presumed.

\begin{proposition}[Rational shifts and cyclotomic order derivatives]
\label{prop:cyclotomic}
Let $1\le r\le q$ and $\omega=e^{2\pi i/q}$. Then
\begin{equation}\label{eq:rational-L}
 \mathcal L_{r/q}(s;t)=q^{s-1}\sum_{h=0}^{q-1}
               \omega^{-rh}\Li_s(\omega^h e^{-t/q}).
\end{equation}
For the ordered depth-two heat sum
\begin{equation}\label{eq:L2-def}
 \mathcal L_a^{(2)}(s_1,s_2;t)=
 \sum_{\nu>\mu\ge0}\frac{e^{-(\nu+a)t}}
                        {(\nu+a)^{s_1}(\mu+a)^{s_2}},
\end{equation}
one has
\begin{align}\label{eq:rational-L2}
 \mathcal L_{r/q}^{(2)}(s_1,s_2;t)
 =q^{s_1+s_2-2}\sum_{h,k=0}^{q-1}\omega^{-r(h+k)}
       \Li_{s_1,s_2}(\omega^h e^{-t/q},\omega^k).
\end{align}
Here $\Li_{s_1,s_2}(z_1,z_2)=\sum_{n>m\ge1}z_1^nz_2^m/(n^{s_1}m^{s_2})$.
Both identities may be differentiated to every fixed order in the spectral
variables and substituted into the ordinary integrals of this article.
\end{proposition}
\begin{proof}
Write $q\nu+r$ for the positive integer in the required residue class.
The root-of-unity filter
$q^{-1}\sum_h\omega^{h(n-r)}$ selects that class. Applying it once
proves \eqref{eq:rational-L}; applying it independently to the two
ordered indices proves \eqref{eq:rational-L2}. For $t>0$ the outer
exponential ensures normal convergence, including every fixed logarithmic
weight. In the resulting finite cyclotomic sums, the endpoint estimates
in Theorem~\ref{thm:C-Mellin} and its depth-two analogue justify the
Mellin substitutions. Derivatives of the prefactor give explicit powers
of $\log q$.
\end{proof}

\section{The complete depth-four normal form}
\label{sec:quartic}
We now answer the predecessor's depth-four question explicitly. To avoid
confusing letters with numerical values, retain $X=Z_0$, $Y=Z_1$,
$Z=Z_2$ for words, and use
\begin{equation}\label{eq:quartic-data}
 x=\gamma_0(a),\quad y=\gamma_1(a),\quad
 z=\zeta(2,a),\quad p=-\zeta'(2,a),\quad r=-\zeta'(3,a).
\end{equation}
All quantities in this section are evaluated at the same positive $a$.
Define
\begin{align}
 A&=\R_a(YXX),&B&=\R_a(XYX),&C&=\R_a(XXY),\label{eq:ABC}\\
 T&=A+B+C,&D&=A-C,&\Theta&=A-2B+C.\label{eq:T-D-Theta}
\end{align}
The final orientation coordinate is $\Omega_2=\mathcal C_{0,2}$.

\begin{proposition}[Permutation components at depth four]
\label{prop:quartic-components}
Among the eight raw coefficients entering a depth-four finite part,
five independent combinations are in the single-letter algebra.
The remaining quotient is
\begin{equation}\label{eq:quartic-representation}
 \operatorname{Std}(S_3)\ \oplus\ \operatorname{sgn}(S_2).
\end{equation}
Here the two groups act on their respective length-three and length-two
word layers; this is not an assertion of a common $S_4$ action preserving
each prefix denominator. The coordinates $(D,\Theta)$ form a rational
basis of the standard component, and $\Omega_2$ spans the sign component.
The symmetric triple value is
\begin{equation}\label{eq:T-evaluation}
 T=yC_2-xp+r
   =\gamma_1(a)\frac{\gamma_0(a)^2-\zeta(2,a)}2
        +\gamma_0(a)\zeta'(2,a)-\zeta'(3,a).
\end{equation}
\end{proposition}
\begin{proof}
The four layers $(j,r)=(1,3),(2,2),(3,1),(4,0)$ contain respectively
$1,3,3,1$ words. The length-two layer consists of $ZX,YY,XZ$:
$ZX+XZ$ and $YY$ are symmetric, and $XZ-ZX$ is the sign component.
The length-three layer is the permutation representation on $YXX,XYX,XXY$.
Its sum is invariant; its sum-zero quotient is the standard two-dimensional
representation. The other layers are invariant single orbits. This proves
the decomposition and the count $1+2+1+1=5$ of symmetric directions.
The change of coordinates in \eqref{eq:T-D-Theta} is invertible, with
\begin{equation}\label{eq:ABC-inverse}
 A=\frac{2T+\Theta}{6}+\frac D2,\quad
 B=\frac{T-\Theta}{3},\quad
 C=\frac{2T+\Theta}{6}-\frac D2.
\end{equation}
Finally the stuffle product of $Y$ with $XX$ is
$YXX+XYX+XXY+(2,1)X+X(2,1)$. The last pair sums to $xp-r$,
while $\R_a(XX)=C_2$. This gives \eqref{eq:T-evaluation}.
\end{proof}

Define the quadratic form
\begin{align}\label{eq:quadratic-Q}
 Q_2(u,v;a)={}&\frac{u^2+v^2}{4}
                 \bigl(x\gamma_2(a)-\zeta''(2,a)\bigr)
             +\frac{v^2-u^2}{4}\Omega_2(a)\notag\\
 &+\frac{uv}{2}\bigl(y^2-\zeta''(2,a)\bigr).
\end{align}

\begin{theorem}[Every ordered quartic finite part]\label{thm:quartic}
For every $a>0$ and every ray with nonzero prefix sums,
\begin{align}\label{eq:quartic-master}
 F_4(c_1,c_2,c_3,c_4;a)={}&C_4
 -\frac1{c_1}\left\{\frac{c_2+c_3+c_4}{3}T
                 +\frac{c_2-c_4}{2}D
                 +\frac{c_2-2c_3+c_4}{6}\Theta\right\}\notag\\
 &+\frac{Q_2(c_3,c_4;a)}{c_1(c_1+c_2)}
 -\frac{c_4^3\gamma_3(a)}{6c_1(c_1+c_2)(c_1+c_2+c_3)}.
\end{align}
Here $C_4$ and $T$ are the single-letter expressions
\eqref{eq:C4} and \eqref{eq:T-evaluation}; the three ordered coordinates
are exactly $D,\Theta,\Omega_2$.
\end{theorem}
\begin{proof}
At $d=4$, formula \eqref{eq:ray} is
\begin{align*}
 F_4={}&\R_a(X^4)
 -\frac{c_2\R_a(YXX)+c_3\R_a(XYX)+c_4\R_a(XXY)}{c_1}\\
 &+\frac{c_3^2\R_a(ZX)/2+c_3c_4\R_a(YY)+c_4^2\R_a(XZ)/2}
          {c_1(c_1+c_2)}
 -\frac{c_4^3\gamma_3(a)}{6c_1(c_1+c_2)(c_1+c_2+c_3)}.
\end{align*}
Use \eqref{eq:ABC-inverse} in the linear term. In the quadratic term,
\eqref{eq:double-normal} gives
\[
 \R_a(ZX)=\frac{x\gamma_2-\zeta''(2,a)-\Omega_2}{2},\quad
 \R_a(XZ)=\frac{x\gamma_2-\zeta''(2,a)+\Omega_2}{2},\quad
 \R_a(YY)=\frac{y^2-\zeta''(2,a)}2.
\]
The result is \eqref{eq:quadratic-Q}, while the constant word is
$C_4$ by \eqref{eq:C-gen}.
\end{proof}

\subsection{Three positive rays and an exact inverse}
Let $K_4(\bbc;a)$ denote the right side of \eqref{eq:quartic-master}
with $D,\Theta,\Omega_2$ set to zero, leaving the explicitly known
single-letter part. Set
\begin{align*}
 \Delta_1&=F_4(1,1,1,2;a)-K_4(1,1,1,2;a),\\
 \Delta_2&=F_4(1,1,2,3;a)-K_4(1,1,2,3;a),\\
 \Delta_3&=F_4(1,2,3,4;a)-K_4(1,2,3,4;a).
\end{align*}

\begin{corollary}[Positive-ray reconstruction]\label{cor:three-rays}
The three orientation coordinates satisfy the exact identities
\begin{equation}\label{eq:three-rays}
 \begin{pmatrix}\Delta_1\\\Delta_2\\\Delta_3\end{pmatrix}
 =\begin{pmatrix}1/2&-1/6&3/8\\1&0&5/8\\1&0&7/12\end{pmatrix}
 \begin{pmatrix}D\\\Theta\\\Omega_2\end{pmatrix},
 \qquad \det=-\frac1{144}.
\end{equation}
In particular,
\begin{equation}\label{eq:ray-inverse}
 \boxed{\quad
 D=-14\Delta_2+15\Delta_3,\quad
 \Theta=-6\Delta_1+12\Delta_2-9\Delta_3,\quad
 \Omega_2=24\Delta_2-24\Delta_3.\quad}
\end{equation}
\end{corollary}
\begin{proof}
In \eqref{eq:quartic-master}, the coefficients of $(D,\Theta,\Omega_2)$
are
\[
 \left(-\frac{c_2-c_4}{2c_1},\;
       -\frac{c_2-2c_3+c_4}{6c_1},\;
        \frac{c_4^2-c_3^2}{4c_1(c_1+c_2)}\right).
\]
Substitution of the three rays gives the displayed matrix. Rational
Gaussian elimination gives its determinant and inverse.
\end{proof}

This corollary provides identities between genuinely ordered Laurent
coefficients and convergent Stieltjes coordinates. No extrapolation from
a diagonal or a fully symmetrized calculation would determine all three.

\subsection{All equal inner slopes and a diagonal check}
For $c+rt\ne0$ at the relevant prefixes,
\begin{align}\label{eq:quartic-closed}
 F_4(c,t,t,t;a)={}&C_4-\frac tcT
 +\frac{t^2}{2c(c+t)}
       \bigl(x\gamma_2(a)+y^2-2\zeta''(2,a)\bigr)\notag\\
 &-\frac{t^3\gamma_3(a)}{6c(c+t)(c+2t)}.
\end{align}
This has no ordered coordinate at all, in agreement with
Theorem~\ref{thm:closure}. In particular,
\begin{equation}\label{eq:quartic-diagonal}
 F_4(1,1,1,1;a)
 =C_4-T+\frac{x\gamma_2(a)+y^2-2\zeta''(2,a)}4
                        -\frac{\gamma_3(a)}{36}.
\end{equation}
An independent exact check is available from the elementary-symmetric
identity
\begin{align}\label{eq:newton4}
 Z_4(s,s,s,s;a)=\frac1{24}\bigl\{&\zeta(s,a)^4
 -6\zeta(s,a)^2\zeta(2s,a)+3\zeta(2s,a)^2\notag\\
 &+8\zeta(s,a)\zeta(3s,a)-6\zeta(4s,a)\bigr\}.
\end{align}
Substitute $s=1+\varepsilon$ and \eqref{eq:stieltjes} and extract the
constant coefficient. The result is exactly \eqref{eq:quartic-diagonal};
the source code performs this as a symbolic rational polynomial identity.

\section{A Gamma-weighted differential identity and antiderivative}
\label{sec:gamma-bridge}
The first standard-component coordinate $D$ has a useful relation to the
first double commutator and to ordinarily convergent generalized harmonic
sums. Define $Q=(2,0)$ and $P=(2,1)$, and put
\begin{equation}\label{eq:H-E-def}
 H(a)=\R_a(QX),\qquad
 E(a)=\R_a(PX)-\R_a(QY).
\end{equation}
For $x_\nu=\nu+a$, $L_\nu=\log x_\nu$, let
\begin{equation}\label{eq:HN-JN}
 h_\nu(a)=\sum_{\mu=0}^{\nu-1}\frac1{\mu+a},\qquad
 j_\nu(a)=\sum_{\mu=0}^{\nu-1}\frac{\log(\mu+a)}{\mu+a}.
\end{equation}
Then
\begin{equation}\label{eq:H-E-sums}
 H(a)=\sum_{\nu\ge0}\frac{h_\nu(a)}{x_\nu^2},\qquad
 E(a)=\sum_{\nu\ge0}\frac{L_\nu h_\nu(a)-j_\nu(a)}{x_\nu^2}.
\end{equation}
These sums converge absolutely and locally uniformly, including the
fixed parameter derivatives used below.

\begin{lemma}[The merged-letter contribution]\label{lem:D-merged}
With the notation \eqref{eq:quartic-data},
\begin{equation}\label{eq:D-merged}
 D=-\frac{x\Omega}{2}+\frac{xp-yz}{2}-E.
\end{equation}
\end{lemma}
\begin{proof}
For transparency, the three finite quasi-shuffle rows needed are
\begin{align*}
 Y*XX&=YXX+XYX+XXY+PX+XP,\\
 X*YX&=2YXX+XYX+PX+YQ,\\
 X*XY&=XYX+2XXY+QY+XP.
\end{align*}
Subtract the last row from the middle row and evaluate. Its left side
is $x(\R_a(YX)-\R_a(XY))=-x\Omega$. On the right,
$PX-QY$ contributes $E$, and
\[
 \R_a(YQ)-\R_a(XP)
 =(yz-r-\R_a(QY))-(xp-r-\R_a(PX))=yz-xp+E.
\]
The triple part contributes $2D$. Rearrangement proves
\eqref{eq:D-merged}.
\end{proof}

This calculation explains why keeping only the highest-length part of a
shuffle relation loses information. The merged letters contribute $E$;
that term does not disappear simply because the top-depth permutation
component is two-dimensional.

\begin{theorem}[Gamma-weighted bridge]\label{thm:gamma-bridge}
For every $a>0$,
\begin{equation}\label{eq:Omega-derivative}
 \Omega'(a)=x(z-p)+yz-\zeta(3,a)-2H(a)+2E(a),
\end{equation}
and consequently
\begin{equation}\label{eq:Gamma-bridge}
 \boxed{\displaystyle
 2D(a)=\gamma_0(a)\zeta(2,a)-\zeta(3,a)-2H(a)
       -\Gamma(a)\frac{\mathrm d}{\mathrm da}
                          \left(\frac{\Omega(a)}{\Gamma(a)}\right).}
\end{equation}
For any positive endpoints $\alpha,\beta$ this gives the exact
antiderivative identity
\begin{align}\label{eq:Gamma-antiderivative}
 \int_\alpha^\beta
 \frac{\gamma_0(a)\zeta(2,a)-\zeta(3,a)-2H(a)-2D(a)}{\Gamma(a)}\dd a
 =\frac{\Omega(\beta)}{\Gamma(\beta)}
  -\frac{\Omega(\alpha)}{\Gamma(\alpha)}.
\end{align}
\end{theorem}
\begin{proof}
The letter derivation gives $\mathcal DX=-Q$ and
$\mathcal DY=-P+Q$. Thus
\begin{align*}
 \Omega'&=-\R_a(QY)-\R_a(XP)+\R_a(XQ)
               +\R_a(PX)-\R_a(QX)+\R_a(YQ).
\end{align*}
Use
\[
 \R_a(XP)=xp-r-\R_a(PX),\quad
 \R_a(YQ)=yz-r-\R_a(QY),\quad
 \R_a(XQ)=xz-\zeta(3,a)-H.
\]
The result is \eqref{eq:Omega-derivative}. Eliminate $E$ with
\eqref{eq:D-merged} to obtain
\[
 2D=xz-\zeta(3,a)-2H-x\Omega-\Omega'.
\]
Since $x=-\psi(a)$,
$\Gamma(a)(\Omega/\Gamma)'=\Omega'+x\Omega$. This proves
\eqref{eq:Gamma-bridge}. Local uniform convergence and the parameter
transport proposition ensure continuous differentiability on every compact
positive interval. The fundamental theorem of calculus yields
\eqref{eq:Gamma-antiderivative}.
\end{proof}

At $a=1$ the classical Euler value $H(1)=\zeta(3)$ makes the bridge
particularly simple:
\begin{equation}\label{eq:D-at-one}
 2D(1)=\gamma\zeta(2)-3\zeta(3)-\gamma\Omega(1)-\Omega'(1).
\end{equation}
For a direct proof of the value used here, the generating identity
$\Li_{1,1}(u,1)=\tfrac12\log^2(1-u)$ gives
\[
 H(1)=\frac12\int_0^1\frac{\log^2(1-u)}u\dd u
     =\frac12\int_0^1\frac{\log^2 v}{1-v}\dd v
     =\zeta(3).
\]
The last step is a positive termwise geometric expansion. The equality
$H(a)=\zeta(3,a)$ has not been used, and is not asserted for general $a$.


\subsection{An all-index Gamma--Stieltjes antiderivative family}
The first-order bridge is one member of a uniform identity. For $q\ge1$
define
\begin{equation}\label{eq:Dq-Hm}
 D_q(a)=\R_a(Z_qXX)-\R_a(XXZ_q),\qquad
 H_m(a)=\R_a((2,m)X)=\sum_{\nu\ge0}
             \frac{\log^m(\nu+a)h_\nu(a)}{(\nu+a)^2}.
\end{equation}
Thus $D_1=D$ and $H_0=H$. The $H_m$ are ordinarily convergent
logarithmically weighted generalized harmonic sums.

\begin{theorem}[All-index Gamma bridge and unit-interval evaluation]
\label{thm:all-Gamma}
For $a>0$ and every integer $q\ge1$,
\begin{align}\label{eq:all-Gamma}
 \Gamma(a)\frac{\mathrm d}{\mathrm da}
       \left(\frac{\mathcal C_{0,q}(a)}{\Gamma(a)}\right)
 =q\bigl\{\gamma_0(a)\ell_{2,q-1}(a)
          -\ell_{3,q-1}(a)-2H_{q-1}(a)\bigr\}-2D_q(a).
\end{align}
In particular, writing $\gamma_q=\gamma_q(1)$,
\begin{equation}\label{eq:unit-Stieltjes-integral}
 \boxed{\displaystyle
 \int_1^2
 \frac{q\{\gamma_0(a)\ell_{2,q-1}(a)-\ell_{3,q-1}(a)-2H_{q-1}(a)\}
                  -2D_q(a)}{\Gamma(a)}\dd a=\gamma_q.}
\end{equation}
The one-variable letters in these formulas are explicitly
$\ell_{k,q-1}(a)=(-1)^{q-1}\zeta^{(q-1)}(k,a)$.
\end{theorem}
\begin{proof}
Put $Y_q=Z_q$, $P_q=(2,q)$, and
$E_q=\R_a(P_qX)-\R_a(QY_q)$. Repeating the three stuffle rows in
Lemma~\ref{lem:D-merged}, with $Y_q$ replacing $Y$, gives
\[
 2D_q=-x\mathcal C_{0,q}+x\ell_{2,q}-\gamma_q(a)z-2E_q.
\]
The derivative rule $\mathcal DY_q=-P_q+qP_{q-1}$ gives
\begin{align*}
 \mathcal C_{0,q}'
 &=2E_q-x\ell_{2,q}+\gamma_q(a)z
          +q\{\R_a(XP_{q-1})-\R_a(P_{q-1}X)\}\\
 &=2E_q-x\ell_{2,q}+\gamma_q(a)z
          +q\{x\ell_{2,q-1}-\ell_{3,q-1}-2H_{q-1}\}.
\end{align*}
Eliminate $E_q$ and use $x=-\psi(a)$ to obtain
\eqref{eq:all-Gamma}. This is an identity of finite quasi-shuffle
expressions under the letter derivation; no convergence claim beyond the
already established regularization and parameter transport is needed.

Integrate the derivative of $\mathcal C_{0,q}/\Gamma$ over $[1,2]$.
Since $\Gamma(1)=\Gamma(2)=1$, its endpoint difference is
$\mathcal C_{0,q}(2)-\mathcal C_{0,q}(1)$. At $a=1$,
\eqref{eq:C-shift} and $q\ge1$ give precisely this difference as
$\gamma_q$. The integration is ordinary: every function involved is
smooth on the compact interval.
\end{proof}

The case $q=1$ is the particularly simple identity
\begin{equation}\label{eq:gamma1-integral}
 \int_1^2\frac{\gamma_0(a)\zeta(2,a)-\zeta(3,a)-2H(a)-2D(a)}{\Gamma(a)}\dd a
 =\gamma_1.
\end{equation}
The appearance of a Stieltjes constant here comes from the exact shift
of an ordered commutator, rather than from a singular endpoint or a
formal exchange of divergent integrals.

\section{The second standard component as a curvature integral}
\label{sec:curvature}
The term ``curvature'' here refers only to the second finite difference
$A-2B+C$ among three placements of a logarithmic letter. It does not
refer to differential geometry or an unproved sign property.

\begin{theorem}[Ordinary Mellin formula for the triple curvature]
\label{thm:curvature}
Define
\begin{align}\label{eq:Xi-integral}
 \Xi(a)=\int_0^\infty K(t)\bigl\{
 &2\partial_{s_1}\mathcal L_a^{(2)}(1,1;t)
   -\partial_{s_2}\mathcal L_a^{(2)}(1,1;t)\notag\\
 &-(\log t+\gamma)\mathcal L_a^{(2)}(1,1;t)\bigr\}\dd t.
\end{align}
The integral is ordinarily absolutely convergent, and
\begin{equation}\label{eq:Theta-Xi}
 \boxed{\displaystyle
 \Theta(a)=\Xi(a)+\frac34\gamma_0(a)\gamma_2(a)
                 -\frac12\gamma_1(a)^2-\frac14\zeta''(2,a)
                 -\frac34\Omega_2(a).}
\end{equation}
\end{theorem}
\begin{proof}
Insert \eqref{eq:B-Mellin} into \eqref{eq:tail-A} at depth three:
\[
 A_3(u,v,w)=\frac1{\Gamma(1+u)}\int_0^\infty
                  K(t)t^u\mathcal L_a^{(2)}(1+v,1+w;t)\dd t.
\]
Write $\HH_2(u,v)=\sum h_{mn}u^mv^n$. The divided difference in
\eqref{eq:H-recursion} then has linear part
$h_{20}(u+2v)+h_{11}w$. By \eqref{eq:jet-word},
\[
 \Theta=-\partial_u\HH_3(0)+2\partial_v\HH_3(0)-\partial_w\HH_3(0)
        =\Xi+3h_{20}-h_{11}.
\]
The double normal form yields
\[
 h_{20}=\frac{x\gamma_2-\zeta''(2,a)-\Omega_2}{4},\qquad
 h_{11}=\frac{y^2-\zeta''(2,a)}2.
\]
Substitution gives \eqref{eq:Theta-Xi}.

For endpoint control, the undifferentiated heat sum in
\eqref{eq:L2-def}, and each first order derivative, is bounded near zero
by a fixed power of $1+|\log t|$; split the outer index at $1/t$ and
bound the inner harmonic sums by logarithmic polynomials. Multiplication
by $\log t$ preserves integrability. At infinity the outer exponential
dominates the inner logarithmic growth. These estimates justify the
integral manipulations and show that no regularized endpoint is hidden
in \eqref{eq:Xi-integral}.
\end{proof}

At $a=1$, the heat sum is $\Li_{s_1,s_2}(e^{-t},1)$; at $(1,1)$ it
is $\tfrac12\log^2(1-e^{-t})$. Consequently
\eqref{eq:Theta-Xi} is an explicit identity involving the first order
derivatives of a depth-two polylogarithm, elementary logarithms, Stieltjes
constants, and a Hurwitz-zeta derivative. Proposition
\ref{prop:cyclotomic} supplies its rational-shift cyclotomic forms.

\begin{proposition}[A convergent harmonic--Stieltjes series for $\Xi$]
\label{prop:Xi-series}
With the notation \eqref{eq:HN-JN}, put $x=\nu+a$, $L=\log x$, and
\begin{equation}\label{eq:b01}
 b_0(x)=\gamma_0(x+1)+L,\qquad
 b_1(x)=-\gamma_1(x+1)-\frac{L^2}{2}.
\end{equation}
Then
\begin{equation}\label{eq:Xi-series}
 \Xi(a)=\sum_{\nu=0}^\infty
       \frac{\bigl(-b_1(x)-2L b_0(x)\bigr)h_\nu(a)
                         +b_0(x)j_\nu(a)}x.
\end{equation}
This is absolutely convergent. Its terms are
$O(x^{-2}(1+\log x)^2)$.
\end{proposition}
\begin{proof}
In the normally convergent double series for $A_3$, the undifferentiated
inner sum is $h_\nu$, while differentiation in the inner exponent gives
$-j_\nu$. Also $B(0,x)=b_0(x)$ and $B_u(0,x)=b_1(x)$ by
\eqref{eq:B-derivative}. Therefore
$-A_{3,u}+2A_{3,v}-A_{3,w}$ is exactly
\eqref{eq:Xi-series}. The bound follows from
$b_0=O(x^{-1})$, $b_1=O(x^{-1}(1+\log x))$ and the elementary harmonic
bounds $h_\nu=O(1+\log x)$, $j_\nu=O((1+\log x)^2)$.
\end{proof}

Together, \eqref{eq:Omega2}, \eqref{eq:D-merged},
\eqref{eq:H-E-sums}, and \eqref{eq:Theta-Xi}--\eqref{eq:Xi-series}
are a complete convergent-series realization of the quartic normal form.
Alternatively, the Gamma bridge trades $D$ for $\Omega'$ and $H$,
and the Mellin formulas trade the other coordinates for ordinary
polylogarithmic integrals. These are equivalent exact representations;
none presupposes an unproved reduction to a smaller constant field.

\section{Worked identities and a comparison of representations}
\label{sec:examples}
The results are more informative when different exact representations are
kept distinct. A ray finite part, a convergent Stieltjes sum, and an ordinary
Mellin integral may encode the same coordinate, but require different
analytic and algebraic arguments.

\subsection{An explicit difference of positive-ray finite parts}
Let
\[
 F^{(2)}=F_4(1,1,2,3;a),\qquad F^{(3)}=F_4(1,2,3,4;a).
\]
Eliminating the known parts in \eqref{eq:ray-inverse} gives the exact
identity
\begin{align}\label{eq:explicit-ray-difference}
 \Omega_2(a)={}&24\bigl(F^{(2)}-F^{(3)}\bigr)-24T
 +11\gamma_0(a)\gamma_2(a)+12\gamma_1(a)^2\notag\\
 &-23\zeta''(2,a)-\frac{13}{18}\gamma_3(a).
\end{align}
For verification by hand, the two single-letter parts differ by
\[
 K_4(1,1,2,3;a)-K_4(1,2,3,4;a)
 =T-\frac{11}{24}\bigl(x\gamma_2-\zeta''(2,a)\bigr)
    -\frac12\bigl(y^2-\zeta''(2,a)\bigr)+\frac{13}{432}\gamma_3.
\]
Multiplying by $24$ proves \eqref{eq:explicit-ray-difference}. Inserting
\eqref{eq:Omega2} therefore yields a further convergent identity:
\begin{align}\label{eq:ray-Stieltjes-series}
 \sum_{\nu\ge0}\frac{L^2\gamma_0(x)-\gamma_2(x)+2L^3/3}{x}
 ={}&24\bigl(F^{(2)}-F^{(3)}\bigr)-24T
       +11\gamma_0(a)\gamma_2(a)+12\gamma_1(a)^2\notag\\
 &-23\zeta''(2,a)-\frac1{18}\gamma_3(a),\qquad x=\nu+a,\quad L=\log x.
\end{align}
This explicitly isolates the length-two sign coordinate from two rays
whose length-three orientation coefficients cancel after subtraction. The equality
is exact even when the individual ordered finite parts do not reduce to
single-letter constants.

\subsection{Every ordered double logarithmic moment}
Combining Theorem~\ref{thm:commutator} with
Corollary~\ref{cor:double-normal} gives, without new notation,
\begin{align}\label{eq:all-double-explicit}
 \R_a(Z_pZ_q)={}&\frac12\gamma_p(a)\gamma_q(a)
 -\frac{(-1)^{p+q}}2\zeta^{(p+q)}(2,a)
 +\frac{\kappa_{p,q}}2\gamma_{p+q+1}(a)\notag\\
 &+\frac12\sum_{\nu\ge0}
 \frac{\log^q x\,\gamma_p(x)-\log^p x\,\gamma_q(x)
                  -\kappa_{p,q}\log^{p+q+1}x}{x}.
\end{align}
Here $x=\nu+a$. For $p=q$, the final two terms vanish. For unequal indices, the sum has
an $x^{-3}$ leading tail after the cancellations built into the formula.
The equality defines neither side by the other: the left side is the
harmonic polynomial constant of the finite nested sum, while the right
side uses ordinary one-variable Stieltjes functions in an absolutely
convergent series.

\subsection{Exact finite shifts}
Iterating \eqref{eq:C-shift}, for any positive integer $N$,
\begin{align}\label{eq:finite-shift}
 \mathcal C_{p,q}(a+N)=\mathcal C_{p,q}(a)
 -\sum_{\nu=0}^{N-1}
 \frac{\log^q(\nu+a)\gamma_p(\nu+a)
       -\log^p(\nu+a)\gamma_q(\nu+a)}{\nu+a}.
\end{align}
The one-letter case of \eqref{eq:word-shift} expresses every
$\gamma_m(\nu+a)$ here as $\gamma_m(a)-H_\nu^{(1,m)}(a)$.
Thus all positive integer translates of the commutator are determined by
one base value and finite generalized harmonic sums. At $a=1$ and $q>0$,
this specializes to
\[
 \mathcal C_{0,q}(2)-\mathcal C_{0,q}(1)=\gamma_q,
\]
which is exactly the endpoint mechanism behind
\eqref{eq:unit-Stieltjes-integral}.

\section{Audit findings, verification, and integration}
\label{sec:audit}
\subsection{What is corrected, sharpened, and left open}
No new mathematical error was established in the inspected predecessor's
first-commutator formula. In particular, its term $-\gamma_2(a)/2$ is
correct. The present all-index formula shows why replacing the contact
term by zero, or extrapolating that first factor unchanged to all indices,
would be wrong: the correct coefficient is
$\kappa_{p,q}=1/(q+1)-1/(p+1)$. This is a safeguard for generalization,
not an allegation that the predecessor contained that error.

The predecessor's depth-four organization question can now be replaced
by Proposition~\ref{prop:quartic-components} and Theorem
\ref{thm:quartic}. Its broader all-depth version has the exact linear
answer in Theorem~\ref{thm:rank}. The words ``independent coordinate''
must remain qualified by the specified formal quotient. In particular,
Lemma~\ref{lem:D-merged} shows that a highest-length or Lyndon-only count
cannot be substituted for the complete normal form without tracking the
merged-letter terms and products of earlier orientation coordinates.

These results do not settle the arithmetic reduction of $\Omega(1)$,
$\Omega(1/2)$, or $\Omega(1/3)$ into a smaller named algebra. They do
not prove the Gaussian $S_6$ relation. The inspected research-programme
chapter records an exact rejection of one particular frozen weight-nine
$S_8$ vector; that rejection is not an independence theorem and does not
reject every possible relation in its basket \cite{repo-discovery}.
The present work does not revisit that certificate. It also makes no
claim of proof-assistant formalization.

\subsection{Exact finite certificates}
The script \src{code/verify_exact.py} uses integer word coefficients
and rational symbolic arithmetic. Its saved output records:
\begin{itemize}[itemsep=2pt]
\item 146 complete symmetric-sum expansions, including repeated logarithmic
indices and every partition-lattice merge in each tested instance;
\item 81 exact contact-coefficient calculations for $0\le p,q\le8$,
and eight full word-algebra certificates of the all-index Gamma bridge;
\item the rank table through $d=12$, the three displayed quartic stuffle
rows, the quartic linear change of coordinates, its positive-ray matrix
and inverse, and the independent diagonal Newton constant-term identity.
\end{itemize}
The all-depth rank and analytic identities are not inferred from these
finite checks. Their proofs are given above. The checks make the finite
algebra and implementation replayable and help guard against signs,
factorials, repeated-index multiplicities, and missing merged letters.

\subsection{Numerical regressions and their limitations}
The script \src{code/verify_numeric.py} performs 83 floating-point
regressions at 80 decimal working digits, using $a=1/2,1,3/2$.
It compares two cutoff/asymptotic settings, $(N,M)=(48,32)$ and $(64,40)$;
checks single letters against separate \texttt{mpmath} Stieltjes and
Hurwitz-zeta calls; compares raw harmonic constants with the convergent
commutator and curvature series; and compares five quartic rays at each
shift with the explicit normal form. A five-point real-parameter derivative
provides a separate check of the Gamma-weighted differential bridge.

The maximum observed residual among the 82 ordinary regressions was
approximately $1.462\times10^{-45}$, dominated by the comparison of
truncation settings. The five-point derivative check had residual
approximately $3.477\times10^{-20}$ with step $10^{-5}$. These are
observed floating-point differences, not proven error bounds. The two
nested representations share the Euler--Maclaurin engine, and this shared
dependency is recorded explicitly in the data. Neither tiny residuals nor
stability under a larger cutoff is used as a proof of equality.

Selected diagnostic values, deliberately rounded far below the working
precision, are:
\begin{center}\small
\begin{tabular}{r r r r r}
\toprule
$a$&$\Omega(a)$&$\Omega_2(a)$&$D(a)$&$\Theta(a)$\\
\midrule
$1/2$&$0.013640874271$&$-0.005542291305$&$-2.652698961380$&$1.490647980568$\\
$1$&$0.093995347735$&$0.041077482932$&$-1.272799745634$&$0.210840232622$\\
$3/2$&$0.028724389737$&$0.045438039415$&$-0.951236020795$&$-0.095917401069$\\
\bottomrule
\end{tabular}
\end{center}
The package preserves longer values, every residual, the settings, and
whether each tolerance check passed. These decimals are diagnostics only;
no sign or arithmetic-independence theorem is inferred from them.

\subsection{Repository-ready contents}
The delivered directory contains this article and PDF; Python verification
code; JSON and text outputs; a provenance manifest; build instructions;
and a short manuscript insertion fragment with namespaced labels. A
suggested report location is
\begin{center}\small\ttfamily
Analysis/Polylogarithms/docs/reports/ordered-jet-normal-forms/.
\end{center}
The fragment is intended for the integration/differentiation part of the
manuscript, with a cross-reference from the research-programme chapter.
Its status note replaces the specified formal depth-four question, while
retaining the arithmetic evaluation questions. No remote repository file
was modified during this work.

\section{Further research questions}
\label{sec:future}
The following problems are deliberately more specific than asking for
additional high-precision fits.

\paragraph{1. Evaluation kernels at rational shifts.}
For $a=1,1/2,1/3$, fix an explicit comparison algebra generated by selected
Gamma values, ordinary zeta derivatives, Dirichlet-$L$ derivatives, and
cyclotomic polylogarithms. Determine which classes from
Theorem~\ref{thm:rank} vanish after numerical evaluation modulo that
algebra. The first concrete targets are $\Omega_2(a)$ and the combination
$\Xi(a)-3\Omega_2(a)/4$. The formal rank is not an obstruction to
additional analytic reductions.

\paragraph{2. Linear quotients versus scalar extension.}
Compute the relevant quotient after products of previously known ordered
coordinates with single-letter data are allowed as background. At depth
four, \eqref{eq:D-merged} provides the first exact transition between
these two questions. A general algorithm should track lower-length merged
words, not just the top-length Lyndon terms. Its output may be a filtered
module presentation rather than the vector-space count $N_d$.

\paragraph{3. Adding convergent shuffle and functional equations.}
The single-letter-closed locus in Theorem~\ref{thm:closure} is complete
for the stipulated stuffle vocabulary. Determine its enlargement when
specific convergent iterated-integral shuffle identities, reflection
identities, or functional equations are added. The target is a new exact
presentation and a replayable certificate, not an inference from failed
integer-relation searches.

\paragraph{4. Colored commutators and conductors.}
Extend the decorated alphabet by root-of-unity colors. Merging then
multiplies colors while adding both exponents. Derive the corresponding
contact terms, character projectors, and ordinary Mellin formulas, including
what changes when the color is nontrivial and the one-variable pole is
absent. The rational-shift formulas of Proposition~\ref{prop:cyclotomic}
provide concrete consistency checks.

\paragraph{5. Nonlinear and tangential finite parts.}
For a fixed analytic curve through the all-one point, the required jet
orders depend on its orders of contact with the prefix divisors. Apply
the orbit-sum theorem to the resulting finite set of bidegrees and derive
an exact orientation count in terms of those contact orders. Unlike the
transverse formula, the set need not lie on the single diagonal $j+r=d$.
A coefficient-identifiability theorem must then account for constraints
among the curve's Taylor coefficients.

\paragraph{6. Eliminating higher ordered terms from the Gamma integrals.}
Formula \eqref{eq:unit-Stieltjes-integral} gives an ordinary integral for
every $\gamma_q$. Seek additional exact identities that eliminate $D_q$
in favor of lower-depth functions or classical polylogarithm order
derivatives. It is essential to specify which function class is permitted;
simply replacing $D_q$ by the derivative in \eqref{eq:all-Gamma} would
only restate the fundamental theorem of calculus.

\paragraph{7. Rational ray designs at higher depth.}
Theorem~\ref{thm:identifiable} proves existence of rational positive sample
rays. Find explicit structured families with transparent inverse matrices
for $d=5,6$, using the nine and twenty-one orientation coordinates,
respectively. The aim is symbolic recoverability and simple exact
coefficients, not a numerical conditioning claim without an accompanying
analysis.

\paragraph{8. Formal verification of the algebraic--analytic interface.}
A first proof-assistant target is the finite word algebra: merge,
quasi-shuffle, partition-lattice symmetrization, the orbit quotient, and
the ray matrix. The analytic target can then isolate normal convergence
of the outer-tail series, the uniqueness of the cutoff polynomial, and
parameter differentiation. This separation avoids treating a successful
finite symbolic script as a formal proof of an analytic identity.

\appendix
\section{The numerical harmonic constant-term evaluator}
\label{app:em}
The evaluator represents a truncated asymptotic expansion in
$x=N+a$ as a dictionary of coefficients of $x^{-p}\log^q x$.
Finite nested sums are evaluated by the exact ordering recurrence
\[
 W_{(k,m)v}(n+1)=W_{(k,m)v}(n)+
                  (n+a)^{-k}\log^m(n+a)W_v(n),\qquad W_\varnothing=1.
\]
For a prospective increment expansion $f(x)$, a formal discrete primitive
is
\begin{equation}\label{eq:EM-primitive}
 \mathfrak E f(x)=\int f(x)\dd x-\frac12f(x)
       +\sum_{r\ge1}\frac{B_{2r}}{(2r)!}f^{(2r-1)}(x).
\end{equation}
Every integration constant is set to zero in the logarithmic polynomial.
For a monomial, the primitives are
\begin{align*}
 \int x^{-1}\log^q x\dd x&=\frac{\log^{q+1}x}{q+1},\\
 \int x^{-p}\log^q x\dd x
 &=-x^{1-p}\sum_{h=0}^q
       \frac{q!\,\log^{q-h}x}{(q-h)!(p-1)^{h+1}},\qquad p>1.
\end{align*}
All derivatives and products of the finite expansions are elementary.
The constant is estimated recursively as the finite sum at $N$ minus
its nonconstant asymptotic primitive at $N+a$. Convergent outer sums use
the same primitive to approximate their tails. For the commutator series,
exactly cancelling nonsummable polynomial terms are removed symbolically
by their exponent pattern before the numerical tail is formed.

This algorithm uses an \emph{asymptotic} expansion. Increasing its order
indefinitely at fixed cutoff is not asserted to converge. The recorded
tests increase both cutoff and order, compare with separate single-variable
calls, and retain all residuals. No interval enclosure or optimal-truncation
bound is implemented. The mathematical proofs in the body of the article
do not depend on the numerical evaluator.

\section{Finite normal-form and ray-recovery algorithms}
\label{app:algorithms}
For a given depth $d$, enumerate all $j=1,\ldots,d$ and weak compositions
$\bbm$ of $d-j$ into $j$ parts. Group the tuples by their sorted version.
For each group, generate the labelled symmetric-sum certificate from
\eqref{eq:symmetric-sum}, divide by the repeated-entry multiplicity, and
apply \eqref{eq:normal-algorithm}. Substitute the monomial coefficient
from \eqref{eq:ray}. All operations are finite rational additions,
multiplications, and coefficient collections.

For a practical positive-ray search, enumerate positive integer vectors
in a fixed order, omit duplicates under common rescaling, and append a
row whenever exact rational elimination increases the orientation-matrix
rank. Theorem~\ref{thm:identifiable} proves that some finite stage attains
rank $N_d$: rational positive rays can be scaled to positive integer ones,
and a common slope scaling leaves \eqref{eq:ray} unchanged. This is a
termination theorem for a specified algebraic enumeration. It is not a
claim that the first full-rank matrix has small entries or is numerically
well conditioned.

The certificate checker expands quasi-shuffle products recursively via
\[
 (av)*(bw)=a\bigl(v*(bw)\bigr)+b\bigl((av)*w\bigr)
                  +(a\circ b)(v*w).
\]
The empty word is the multiplicative unit. Word coefficients are integers,
so a vanishing residual is checked by exact coefficient cancellation.
The explicit $3\times3$ matrix in \eqref{eq:three-rays} and the formal
Gamma bridges are examples of certificates produced without numerical
constants of any kind.

\section{A compact theorem dependency map}
\label{app:dependencies}
The analytic foundation consists of the harmonic constant-term character,
the outer-tail Mellin representation, and the compatible germ. The
orientation rank uses only the universal quasi-shuffle algebra and the
partition-lattice symmetric sum; it does not require analytic continuation.
Ray identifiability uses only rational functions of the slopes. The ray
finite-part interpretation then joins these two foundations.

The commutator series is obtained from the depth-two outer-tail recursion.
The quartic formula is a finite substitution into the ray coefficient
formula. The Gamma bridge uses the parameter derivative of words and
three finite stuffle rows. The curvature identity uses the depth-three
outer-tail recursion. Thus none of the new analytic identities is justified
by the numerical regression table, and the exact rank theorem does not
assume any unproved arithmetic property of zeta or Stieltjes values.

\clearpage
\begin{thebibliography}{99}\raggedright
\bibitem{ordered}
\emph{Ordered Hurwitz Germs at Every Depth: Harmonic renormalization,
directional Stieltjes constants, and polylogarithmic Mellin identities}.
ProveIt research continuation, 10 October 2026. Inspected complete TeX
\texttt{article(20261011-014723).tex}; particularly Sections 4--7 and
further-research question 2. The visible incoming inventory includes
\texttt{ProveIt\_Ordered\_Hurwitz\_Germs\_2026-10-10.zip}. Source hashes
and inspection scope are recorded in the accompanying provenance manifest.

\bibitem{repo-manuscript}
V. Reshetnikov, \emph{ProveIt: Polylogarithms manuscript},
\nolinkurl{Analysis/Polylogarithms/docs/manuscript/}, inspected 10 October 2026.
\url{https://github.com/VladimirReshetnikov/ProveIt/tree/main/Analysis/Polylogarithms/docs/manuscript}.

\bibitem{repo-discovery}
V. Reshetnikov, \emph{Experimental discovery and the route to proof},
ProveIt manuscript, \texttt{chapters/10-discovery.tex}. Inspected Git blob
\texttt{207c719ae6e7fb1639401f449216ae9b695afd40}.
\url{https://github.com/VladimirReshetnikov/ProveIt/blob/main/Analysis/Polylogarithms/docs/manuscript/chapters/10-discovery.tex}.

\bibitem{hoffman}
M. E. Hoffman, \emph{Quasi-shuffle products}, Journal of Algebraic
Combinatorics \textbf{11} (2000), 49--68.
\url{https://arxiv.org/abs/math/9907173}.

\bibitem{mow}
K. Matsumoto, T. Onozuka, and I. Wakabayashi,
\emph{Laurent series expansions of multiple zeta-functions of Euler--Zagier
type at integer points}, preprint, arXiv:1601.05918 (2016).
\url{https://arxiv.org/abs/1601.05918}.

\bibitem{dlmf-zeta}
NIST Digital Library of Mathematical Functions, Section 25.11,
\emph{Hurwitz Zeta Function}; definitions, recurrence, derivatives,
and Euler--Maclaurin representation.
\url{https://dlmf.nist.gov/25.11}.

\bibitem{dlmf-polylog}
NIST Digital Library of Mathematical Functions, Section 25.12,
\emph{Polylogarithms}.
\url{https://dlmf.nist.gov/25.12}.

\bibitem{dlmf-lerch}
NIST Digital Library of Mathematical Functions, Section 25.14,
\emph{Lerch's Transcendent}.
\url{https://dlmf.nist.gov/25.14}.

\bibitem{dlmf-gamma}
NIST Digital Library of Mathematical Functions, Sections 5.7 and 5.15,
\emph{Gamma-function series expansions} and \emph{Polygamma Functions}.
\url{https://dlmf.nist.gov/5.7}; \url{https://dlmf.nist.gov/5.15}.
\end{thebibliography}
\end{document}
